% ============================================================
\chapter{Large-Scale Training --- Distributed Training}
\label{ch:training}
\index{distributed training}
% ============================================================

\section{Why Distributed Training?}
\label{sec:why-distributed}

Modern deep learning models (LLMs, vision transformers, diffusion models)
have billions of parameters and require terabytes of training data.
A single GPU cannot fit the model or process the data in a reasonable time.
Distributed training splits the workload across multiple devices.

\begin{definition}[Distributed training]
\textbf{Distributed training}\index{distributed training} is the practice
of training a model across multiple GPUs, machines, or accelerators by
partitioning the computation (data, model, or both) and synchronising
the results.
\end{definition}

\section{Data Parallelism vs Model Parallelism}
\label{sec:parallelism-types}
\index{data parallelism}\index{model parallelism}

\begin{center}
\begin{tabular}{@{}lp{5.5cm}p{5.5cm}@{}}
\toprule
& \textbf{Data Parallelism} & \textbf{Model Parallelism} \\
\midrule
\textbf{Principle} & Same model on each device, different data &
  Different parts of the model on different devices \\
\textbf{When} & Model fits in one GPU memory &
  Model too large for one GPU \\
\textbf{Communication} & Gradient synchronisation (AllReduce) &
  Activation exchange between devices \\
\textbf{Scaling} & Near-linear with $N$ GPUs &
  Sub-linear (pipeline bubbles) \\
\textbf{Complexity} & Low & High \\
\textbf{Tools} & DDP, DeepSpeed ZeRO & Megatron-LM, pipeline parallelism \\
\bottomrule
\end{tabular}
\end{center}

\subsection{Data parallelism --- visual overview}

\begin{center}
\begin{tikzpicture}[
  gpu/.style={rectangle, draw, rounded corners, minimum width=2.5cm,
    minimum height=2cm, align=center, font=\small, fill=blue!10},
  data/.style={rectangle, draw, minimum width=1.8cm,
    minimum height=0.6cm, align=center, font=\scriptsize, fill=green!10},
  arr/.style={-{Stealth[length=2mm]}, thick, steelblue},
  lbl/.style={font=\scriptsize, fill=white, inner sep=1pt}
]
  % Data
  \node[data, minimum width=8cm] (batch) {\textbf{Mini-batch} (size $B$)};

  % Split
  \node[data, below left=1.2cm and 1.5cm of batch] (d0)
    {Sub-batch $B/N$};
  \node[data, below=1.2cm of batch] (d1) {Sub-batch $B/N$};
  \node[data, below right=1.2cm and 1.5cm of batch] (d2)
    {Sub-batch $B/N$};

  \draw[arr] (batch.south) -- (d0.north);
  \draw[arr] (batch.south) -- (d1.north);
  \draw[arr] (batch.south) -- (d2.north);

  % GPUs
  \node[gpu, below=0.8cm of d0] (g0) {\textbf{GPU 0}\\Model copy\\$\nabla_0$};
  \node[gpu, below=0.8cm of d1] (g1) {\textbf{GPU 1}\\Model copy\\$\nabla_1$};
  \node[gpu, below=0.8cm of d2] (g2) {\textbf{GPU 2}\\Model copy\\$\nabla_2$};

  \draw[arr] (d0) -- (g0);
  \draw[arr] (d1) -- (g1);
  \draw[arr] (d2) -- (g2);

  % AllReduce
  \node[rectangle, draw, rounded corners, fill=orange!15,
    minimum width=8cm, minimum height=0.8cm, below=1cm of g1,
    font=\small] (ar) {\textbf{AllReduce}: $\bar{\nabla} =
    \frac{1}{N}\sum_{i=0}^{N-1} \nabla_i$};

  \draw[arr, deepred] (g0.south) -- (ar.north);
  \draw[arr, deepred] (g1.south) -- (ar.north);
  \draw[arr, deepred] (g2.south) -- (ar.north);

  % Update
  \node[font=\small, below=0.5cm of ar] {All GPUs update weights with
    $\bar{\nabla}$};
\end{tikzpicture}
\end{center}

\section{Mathematics of Data-Parallel SGD}
\label{sec:dp-math}

Consider $N$ workers, each computing gradients on a sub-batch of size
$b = B/N$ where $B$ is the global batch size.

\begin{keyformulas}[title=\textbf{Data-parallel SGD}]
Each worker $k$ computes:
\[
  \nabla_k = \frac{1}{b} \sum_{i \in \mathcal{B}_k}
    \nabla_\theta \ell(f_\theta(x_i), y_i)
\]
where $\mathcal{B}_k$ is the sub-batch assigned to worker $k$.

The averaged gradient is:
\[
  \bar{\nabla} = \frac{1}{N} \sum_{k=0}^{N-1} \nabla_k
    = \frac{1}{B} \sum_{i \in \mathcal{B}}
    \nabla_\theta \ell(f_\theta(x_i), y_i)
\]
which is identical to computing the gradient on the full batch $\mathcal{B}$.

The parameter update is:
\[
  \theta_{t+1} = \theta_t - \eta \, \bar{\nabla}
\]

\textbf{Linear scaling rule} (Goyal et al., 2017): when multiplying the
batch size by $N$, multiply the learning rate by $N$ and use a warm-up
period:
\[
  \eta_{\text{distributed}} = N \cdot \eta_{\text{single}}
\]
\end{keyformulas}

\begin{remark}
The AllReduce operation ensures that all workers have identical gradients
(and thus identical model weights) after each step. This is equivalent
to training on a single device with batch size $B = N \cdot b$.
\end{remark}

\subsection{AllReduce communication}

\begin{center}
\begin{tikzpicture}[
  gpu/.style={circle, draw, minimum size=1.2cm, align=center,
    font=\scriptsize, fill=blue!10},
  arr/.style={-{Stealth[length=2mm]}, thick}
]
  % Ring AllReduce
  \node[gpu] (g0) at (0,0) {GPU 0\\$\nabla_0$};
  \node[gpu] (g1) at (3,0) {GPU 1\\$\nabla_1$};
  \node[gpu] (g2) at (3,-2.5) {GPU 2\\$\nabla_2$};
  \node[gpu] (g3) at (0,-2.5) {GPU 3\\$\nabla_3$};

  \draw[arr, steelblue] (g0) -- (g1)
    node[midway, above, font=\scriptsize] {reduce};
  \draw[arr, steelblue] (g1) -- (g2)
    node[midway, right, font=\scriptsize] {reduce};
  \draw[arr, steelblue] (g2) -- (g3)
    node[midway, below, font=\scriptsize] {reduce};
  \draw[arr, steelblue] (g3) -- (g0)
    node[midway, left, font=\scriptsize] {reduce};

  \node[below=0.5cm of g2, font=\small, anchor=north west,
    xshift=-1.5cm] {
    \textbf{Ring AllReduce}: $O\!\left(\frac{2(N-1)}{N} \cdot M\right)$
    communication};
\end{tikzpicture}
\end{center}

\section{PyTorch Distributed Data Parallel (DDP)}
\label{sec:ddp}
\index{DDP}\index{PyTorch!DDP}

\textbf{DistributedDataParallel}\index{DDP} (DDP) is PyTorch's recommended
approach for multi-GPU training. It uses NCCL for GPU-to-GPU communication
and overlaps gradient computation with communication for efficiency.

\begin{codeblock}[title=\textbf{Complete PyTorch DDP training script}]
\begin{minted}{python}
import os
import torch
import torch.nn as nn
import torch.distributed as dist
from torch.nn.parallel import DistributedDataParallel as DDP
from torch.utils.data import DataLoader, DistributedSampler
from torchvision import datasets, transforms

def setup(rank, world_size):
    """Initialise the distributed process group."""
    os.environ["MASTER_ADDR"] = "localhost"
    os.environ["MASTER_PORT"] = "12355"
    dist.init_process_group(
        backend="nccl", rank=rank, world_size=world_size
    )
    torch.cuda.set_device(rank)

def cleanup():
    dist.destroy_process_group()

class SimpleCNN(nn.Module):
    def __init__(self, num_classes=10):
        super().__init__()
        self.features = nn.Sequential(
            nn.Conv2d(1, 32, 3, padding=1), nn.ReLU(),
            nn.MaxPool2d(2),
            nn.Conv2d(32, 64, 3, padding=1), nn.ReLU(),
            nn.MaxPool2d(2),
        )
        self.classifier = nn.Sequential(
            nn.Linear(64 * 7 * 7, 256), nn.ReLU(),
            nn.Linear(256, num_classes),
        )

    def forward(self, x):
        x = self.features(x)
        x = x.view(x.size(0), -1)
        return self.classifier(x)

def train(rank, world_size, epochs=10):
    setup(rank, world_size)

    # Data
    transform = transforms.Compose([
        transforms.ToTensor(),
        transforms.Normalize((0.1307,), (0.3081,)),
    ])
    dataset = datasets.MNIST(
        "data", train=True, download=True, transform=transform
    )
    sampler = DistributedSampler(
        dataset, num_replicas=world_size, rank=rank
    )
    loader = DataLoader(
        dataset, batch_size=64, sampler=sampler, num_workers=4
    )

    # Model
    model = SimpleCNN().to(rank)
    model = DDP(model, device_ids=[rank])

    optimizer = torch.optim.Adam(model.parameters(), lr=1e-3)
    criterion = nn.CrossEntropyLoss()

    # Training loop
    for epoch in range(epochs):
        sampler.set_epoch(epoch)  # Shuffle differently each epoch
        model.train()
        total_loss = 0

        for batch_idx, (data, target) in enumerate(loader):
            data, target = data.to(rank), target.to(rank)

            optimizer.zero_grad()
            output = model(data)
            loss = criterion(output, target)
            loss.backward()
            optimizer.step()

            total_loss += loss.item()

        if rank == 0:
            avg_loss = total_loss / len(loader)
            print(f"Epoch {epoch}: loss = {avg_loss:.4f}")

    # Save model (only on rank 0)
    if rank == 0:
        torch.save(
            model.module.state_dict(), "model_ddp.pt"
        )

    cleanup()

if __name__ == "__main__":
    world_size = torch.cuda.device_count()
    torch.multiprocessing.spawn(
        train, args=(world_size,), nprocs=world_size
    )
\end{minted}
\end{codeblock}

\begin{codeblock}[title=\textbf{Launching DDP training}]
\begin{minted}{bash}
# Single node, multiple GPUs
torchrun --nproc_per_node=4 train_ddp.py

# Multi-node (2 nodes, 4 GPUs each)
# On node 0:
torchrun --nproc_per_node=4 --nnodes=2 --node_rank=0 \
    --master_addr=node0 --master_port=12355 train_ddp.py
# On node 1:
torchrun --nproc_per_node=4 --nnodes=2 --node_rank=1 \
    --master_addr=node0 --master_port=12355 train_ddp.py
\end{minted}
\end{codeblock}

\begin{warning}[title=\textbf{DDP pitfalls}]
\begin{itemize}
  \item Always call \cli{sampler.set\_epoch(epoch)} to ensure proper
    shuffling across epochs.
  \item Save checkpoints only on rank 0 to avoid file corruption.
  \item Access the underlying model via \cli{model.module} (not
    \cli{model}) when saving or inspecting.
  \item Ensure all processes reach the same synchronisation points
    (no conditional code paths that differ across ranks).
\end{itemize}
\end{warning}

\section{Hugging Face Accelerate}
\label{sec:accelerate}
\index{Accelerate}

\textbf{Accelerate}\index{Accelerate} (by Hugging Face) provides a
high-level abstraction over DDP, FSDP, and DeepSpeed. A single training
script works on CPU, single GPU, multi-GPU, and multi-node without code
changes.

\begin{codeblock}[title=\textbf{Training with Accelerate}]
\begin{minted}{python}
from accelerate import Accelerator
import torch
import torch.nn as nn
from torch.utils.data import DataLoader
from torchvision import datasets, transforms

accelerator = Accelerator()

# Model, optimizer, data (standard PyTorch)
model = SimpleCNN()
optimizer = torch.optim.Adam(model.parameters(), lr=1e-3)
criterion = nn.CrossEntropyLoss()

transform = transforms.Compose([
    transforms.ToTensor(),
    transforms.Normalize((0.1307,), (0.3081,)),
])
dataset = datasets.MNIST(
    "data", train=True, download=True, transform=transform
)
loader = DataLoader(dataset, batch_size=64, shuffle=True)

# Prepare for distributed training (one line!)
model, optimizer, loader = accelerator.prepare(
    model, optimizer, loader
)

# Standard training loop (no rank checks needed)
for epoch in range(10):
    model.train()
    for data, target in loader:
        optimizer.zero_grad()
        output = model(data)
        loss = criterion(output, target)
        accelerator.backward(loss)
        optimizer.step()

    if accelerator.is_main_process:
        print(f"Epoch {epoch} complete")

# Save
accelerator.save_model(model, "model_accelerate/")
\end{minted}
\end{codeblock}

\begin{codeblock}[title=\textbf{Accelerate configuration and launch}]
\begin{minted}{bash}
# Interactive configuration
accelerate config
# Generates ~/.cache/huggingface/accelerate/default_config.yaml

# Launch training
accelerate launch train_accelerate.py

# Launch with overrides
accelerate launch --num_processes 4 --mixed_precision fp16 \
    train_accelerate.py
\end{minted}
\end{codeblock}

\section{Mixed Precision Training (AMP)}
\label{sec:amp}
\index{mixed precision}\index{AMP}

\textbf{Mixed precision training}\index{mixed precision} uses both FP16
(or BF16) and FP32 arithmetic. FP16 is used for forward and backward
passes (faster, less memory), while FP32 is used for the master weights
and loss scaling (numerical stability).

\begin{keyformulas}[title=\textbf{Mixed precision benefits}]
\begin{itemize}
  \item \textbf{Memory}: FP16 tensors use $\frac{1}{2}$ the memory of
    FP32 $\Rightarrow$ larger batch sizes or larger models.
  \item \textbf{Speed}: modern GPUs (Ampere, Hopper) have dedicated FP16
    / BF16 tensor cores, providing $2\text{--}8\times$ throughput.
  \item \textbf{Loss scaling}: gradients in FP16 can underflow.
    A \textbf{loss scaler} multiplies the loss by a large factor before
    backpropagation, then divides gradients before the optimizer step.
\end{itemize}
\end{keyformulas}

\begin{codeblock}[title=\textbf{PyTorch AMP example}]
\begin{minted}{python}
import torch
from torch.amp import autocast, GradScaler

model = SimpleCNN().cuda()
optimizer = torch.optim.Adam(model.parameters(), lr=1e-3)
scaler = GradScaler("cuda")

for data, target in train_loader:
    data, target = data.cuda(), target.cuda()

    optimizer.zero_grad()

    # Forward pass in FP16
    with autocast("cuda"):
        output = model(data)
        loss = criterion(output, target)

    # Backward pass with loss scaling
    scaler.scale(loss).backward()
    scaler.step(optimizer)
    scaler.update()
\end{minted}
\end{codeblock}

\begin{remark}
BF16 (bfloat16) has the same exponent range as FP32, making it more
numerically stable than FP16. It is supported on Ampere GPUs and later
(A100, H100). When available, prefer BF16 over FP16 to avoid loss scaling
altogether.
\end{remark}

\section{Hyperparameter Optimisation with Optuna}
\label{sec:optuna}
\index{Optuna}\index{hyperparameter optimisation}

\textbf{Optuna}\index{Optuna} is a hyperparameter optimisation (HPO)
framework that uses Bayesian optimisation (Tree-structured Parzen
Estimators) to efficiently search the hyperparameter space.

\begin{codeblock}[title=\textbf{Optuna hyperparameter search}]
\begin{minted}{python}
import optuna
import torch
import torch.nn as nn
from sklearn.metrics import accuracy_score

def objective(trial):
    # Suggest hyperparameters
    lr = trial.suggest_float("lr", 1e-5, 1e-1, log=True)
    n_layers = trial.suggest_int("n_layers", 1, 4)
    hidden_size = trial.suggest_categorical(
        "hidden_size", [64, 128, 256, 512]
    )
    dropout = trial.suggest_float("dropout", 0.1, 0.5)
    optimizer_name = trial.suggest_categorical(
        "optimizer", ["Adam", "SGD", "AdamW"]
    )

    # Build model dynamically
    layers = []
    in_features = 784
    for i in range(n_layers):
        layers.append(nn.Linear(in_features, hidden_size))
        layers.append(nn.ReLU())
        layers.append(nn.Dropout(dropout))
        in_features = hidden_size
    layers.append(nn.Linear(in_features, 10))
    model = nn.Sequential(*layers).cuda()

    optimizer = getattr(torch.optim, optimizer_name)(
        model.parameters(), lr=lr
    )
    criterion = nn.CrossEntropyLoss()

    # Train
    for epoch in range(10):
        model.train()
        for data, target in train_loader:
            data, target = data.cuda(), target.cuda()
            optimizer.zero_grad()
            loss = criterion(model(data.view(data.size(0), -1)),
                             target)
            loss.backward()
            optimizer.step()

        # Pruning: report intermediate value
        val_acc = evaluate(model, val_loader)
        trial.report(val_acc, epoch)
        if trial.should_prune():
            raise optuna.TrialPruned()

    return val_acc

# Run the study
study = optuna.create_study(
    direction="maximize",
    pruner=optuna.pruners.MedianPruner(),
    storage="sqlite:///optuna.db",
    study_name="mnist-hpo",
)
study.optimize(objective, n_trials=100, timeout=3600)

# Best result
print(f"Best accuracy: {study.best_value:.4f}")
print(f"Best params: {study.best_params}")
\end{minted}
\end{codeblock}

\begin{bestpractice}[title=\textbf{HPO best practices}]
\begin{itemize}
  \item Use \textbf{pruning} (early stopping of bad trials) to save
    compute---Optuna's median pruner is a good default.
  \item Use \textbf{log-scale} for learning rates and weight decay.
  \item Start with a \textbf{broad search}, then narrow around promising
    regions.
  \item Use a \textbf{persistent storage} (\cli{sqlite:///optuna.db})
    so studies can be resumed.
  \item Combine with distributed training: run each trial on a GPU.
\end{itemize}
\end{bestpractice}

\section{Distributed Training Framework Comparison}
\label{sec:distributed-comparison}

\begin{center}
\begin{tabular}{@{}lccc@{}}
\toprule
\textbf{Criterion} & \textbf{DDP} & \textbf{FSDP} & \textbf{DeepSpeed} \\
\midrule
Paradigm & Data parallel & Fully sharded & ZeRO stages 1--3 \\
Memory efficiency & Baseline & High & Very high \\
Model size & $\le$ 1 GPU & $>$ 1 GPU & $>$ 1 GPU \\
Complexity & Low & Medium & Medium \\
Mixed precision & Manual (AMP) & Built-in & Built-in \\
Offloading (CPU/NVMe) & No & No & Yes \\
Framework & PyTorch & PyTorch & PyTorch + config \\
Best for & Standard training & Large models & Very large models \\
\bottomrule
\end{tabular}
\end{center}

\begin{remark}
\textbf{FSDP} (Fully Sharded Data Parallel) shards model parameters,
gradients, and optimizer states across GPUs---similar to DeepSpeed ZeRO
Stage~3 but natively integrated into PyTorch. For most use cases, start
with DDP; move to FSDP or DeepSpeed only when the model does not fit in
a single GPU.
\end{remark}

\section{Best Practices}

\begin{bestpractice}[title=\textbf{Distributed training golden rules}]
\begin{enumerate}
  \item \textbf{Start simple}: single GPU $\to$ DDP $\to$ FSDP/DeepSpeed.
  \item \textbf{Linear scaling rule}: scale learning rate with the number
    of GPUs and use warm-up.
  \item \textbf{Use mixed precision}: always enable AMP (FP16 or BF16)
    for modern GPUs.
  \item \textbf{Profile before optimising}: use PyTorch Profiler or
    \cli{nvidia-smi} to identify bottlenecks.
  \item \textbf{Checkpoint regularly}: save checkpoints every $N$ steps,
    not just at epoch boundaries.
  \item \textbf{Use Accelerate}: for portability across hardware
    configurations.
\end{enumerate}
\end{bestpractice}

\begin{antipattern}[title=\textbf{Common anti-patterns}]
\begin{itemize}
  \item Training in FP32 on modern GPUs (wasting memory and speed)
  \item Grid search instead of Bayesian HPO (exponentially wasteful)
  \item Not using a \cli{DistributedSampler} (data duplication across
    workers)
  \item Saving checkpoints on all ranks (file corruption, wasted I/O)
  \item Ignoring the linear scaling rule (divergence with large batch
    sizes)
\end{itemize}
\end{antipattern}

\section{Mini-project: Distributed Training Pipeline}
\label{sec:mini-project-ch6}

\begin{miniproject}[title=\textbf{Mini-project 6: Multi-GPU training}]
Train an image classifier on CIFAR-10 or Fashion-MNIST:
\begin{enumerate}
  \item Write a single-GPU training script with mixed precision (AMP).
  \item Convert it to DDP using \cli{torchrun}.
  \item Rewrite it using Hugging Face Accelerate and verify identical
    results.
  \item Run an Optuna HPO sweep (at least 20 trials) over learning rate,
    batch size, and architecture (number of layers, hidden size).
  \item Log all experiments to MLflow or W\&B.
  \item Compare wall-clock time: 1~GPU vs 2~GPUs vs 4~GPUs.
  \item (Bonus) Try FSDP or DeepSpeed for a larger model
    (e.g., ResNet-50).
\end{enumerate}
\end{miniproject}

\section{Exercises}
\label{sec:exercises-ch6}

\begin{exercise}[$\bigstar$ --- Mixed precision basics]
Take an existing single-GPU training script and add mixed precision
training using \cli{torch.amp}. Compare:
\begin{enumerate}
  \item Training time per epoch (FP32 vs FP16 vs BF16)
  \item Peak GPU memory usage (\cli{torch.cuda.max\_memory\_allocated()})
  \item Final model accuracy
\end{enumerate}
Report your findings in a table.
\end{exercise}

\begin{exercise}[$\bigstar\bigstar$ --- DDP implementation]
Implement a DDP training script from scratch (without Accelerate) for a
ResNet-18 on CIFAR-10:
\begin{enumerate}
  \item Initialise the process group and wrap the model with DDP
  \item Use \cli{DistributedSampler} and set the epoch correctly
  \item Add gradient clipping and learning rate scheduling
  \item Save checkpoints only on rank 0
  \item Verify that training on 2 GPUs with batch size 32 per GPU produces
    similar results to 1 GPU with batch size 64
\end{enumerate}
\end{exercise}

\begin{exercise}[$\bigstar\bigstar\bigstar$ --- Full-scale training system]
Build a complete training system for a medium-scale model:
\begin{enumerate}
  \item Use Accelerate with automatic mixed precision
  \item Implement Optuna HPO with pruning and a persistent SQLite study
  \item Add DVC data versioning and MLflow experiment tracking
  \item Implement checkpointing with automatic resume on failure
  \item Write a Prefect flow that orchestrates data preparation, HPO,
    and final training with the best hyperparameters
  \item Produce a final report comparing at least 10 configurations
\end{enumerate}
\end{exercise}

\section{Cheatsheet}

\begin{cheatsheet}[title=\textbf{Chapter 6 Summary}]
\begin{tabular}{@{}p{4cm}p{8.5cm}@{}}
\toprule
\textbf{Tool / Concept} & \textbf{Key Points} \\
\midrule
\textbf{Data parallelism} & Same model on each GPU, AllReduce gradients,
  scale LR linearly \\
\textbf{DDP} & \cli{torchrun --nproc\_per\_node=N},
  \cli{DistributedDataParallel}, \cli{DistributedSampler} \\
\textbf{Accelerate} & \cli{accelerator.prepare()},
  \cli{accelerate launch}, hardware-agnostic \\
\textbf{Mixed precision} & \cli{torch.amp.autocast}, \cli{GradScaler},
  BF16 preferred over FP16 \\
\textbf{Optuna} & \cli{study.optimize()}, \cli{trial.suggest\_*()},
  pruning, persistent storage \\
\textbf{FSDP / DeepSpeed} & For models too large for a single GPU;
  shard parameters + gradients + optimizer \\
\bottomrule
\end{tabular}
\end{cheatsheet}
